\section{Position Paper Guidelines}

\lead{Much like most other committees at a MUN conference, we highly recommend that you submit a position paper to the dais to show your understanding of both the committee itself and your allocation.}

In the case of press, your submission will focus primarily on your knowledge of your assigned committee groupings and your understanding of your assigned publications position. With this in mind, your position paper will take the form of a written article focusing on a topic of your choosing that relates to your committee groupings. This article should be written in the style, and with the potential biases of, your assigned publication. For example, if you were representing Fox News, you may choose to write an article focusing on the republican position on the topic of privatisation in space exploration (the topic for this year's House of Representatives). Alongside this guide of what to focus on in your position paper, we also want to give you some more help and guidance on the writing process itself:

\begin{itemize}
  \item Your position paper should be one or two pages at most. As this is the first piece you will be producing for the committee we don’t want you to overwork yourselves before the conference begins. There will be plenty of opportunity for more expansive long-form content during the conference itself and the main goal of this piece is to familiarise us with your writing and research styles and get you feedback quicker.
  \item Please do read some articles from your assigned outlet to familiarise yourselves with their style and biases, but do not re-hash the topic of an already existing article. Instead, try to put a new spin on a pre-existing topic, or take a recent event adjacent to one of your assigned committees and base your position paper around it.
  \begin{itemize}
    \item Please ground your articles in real world events. As tempting as it may be to invent events or incidents for your article, please write with real events at the forefront of your mind.
  \end{itemize}
  \item As these position papers differ significantly from standard MUN committees, a bibliography is not an essential component. However, if you do make use of statistics or directly reference another article in your work, do cite these sources in accordance with the referencing guide for your citation style.
  \item Finally, do not utilise generative AI or large-language models under any circumstance. As you should be aware this is a breach of IEUMUN 2026’s academic integrity guidelines and on top of this prevents us from providing you with constructive feedback. If you are struggling with your work please do ask us, it’s kind of the whole reason we’re here!
\end{itemize}

Aside from these key points, the content, subject, and structure of your position papers is left in your capable hands. Remember, press offers you the rare opportunity to express your creativity in the setting of MUN, so please please please use the position paper as a first chance to explore and express yourselves. While the submission of a position paper is required in order to be eligible for awards, we cannot state enough that this is a low-stress opportunity to test out your ideas, get a hang of writing for press, and receive some vital feedback to make the conference that little bit easier!
